\scanpage{55}
% archive foliation: 54
On dit alors que le morphisme \deleted{$\unsure{f}$} plat $T\to S$ est \uline{homogène}. \inserted{si $\nu(s)$ est loc. constant.} Cette condition est [dans le cas des morphismes \uline{finis} du moins] invariante par composition des morphismes, extension de la base [trivial sur (ii), \unsure{puis sur les définitions}]

\deleted{Corollaire}
\deleted{\uline{Proposition} Soient \deleted{$\unsure{X,Y,T,S}$} $T,S,X,Y$ comme dans (i) de Prop. précédente. Supposons que $X$ soit un schéma en groupes sur $Y$, et $P$ un fibré principal sur $X$}

\uline{Cor. 4} Sous les conditions de (i), supposons $T/S$ fini, \deleted{et que $\unsure{f}$ \unsure{soit|sur} \illegible{1}} \inserted{si $g$ est surjectif}. Alors $f_*(g)$ est \uline{surjectif}

\uline{Proposition} Le foncteur $f_*$ transforme fibrés \inserted{\uline{\unsure{faibles|faiblement}}} principaux sur un $Y/T$ [sous réserve d'existence \inserted{et sous cette restriction}] en fibrés \unsure{faibles|faiblement} principaux.
\margin{\unsure{Première assertion}, \unsure{conséquence triviale des définitions}}
Il transforme fibrés principaux simples en fibrés principaux simples, sous réserve que $T/S$ soit fini.
Il transforme \deleted{fibrés} fibrés principaux à groupe \inserted{et} \deleted{\unsure{fini}} \unsure{simples|simple} \unsure{pour} \illegible{1} en \deleted{\unsure{des}} des
